\section{A 688-bit equational presentation: the truth axiom is derived}
\label{sec:e688}

The strongest alternative formalism among the five inspected incoming
archives is the complete $713$-bit packet $\mathcal E_{713}$, supplied in
the archive \texttt{zfc238-composition-20261019}. It has eight equations,
twelve operation symbols, and the ground query
\[
                 N(d,d)=d.
\]
The packet includes its ordered signature, arities, all equations, and both
endpoints. Its ordinary consistency calibration is independent of the
binary-state constructions. The $946$- and $937$-bit packets in other
archives are earlier controls, so they are not the best baseline for this
comparison.

There is a further reduction within the \emph{same} signature and codec:
the explicit equation forcing $d$ to be Boolean truth follows from the
other equations. Deleting it gives a complete $688$-bit packet with seven
equations. The proof below establishes equality of the two model classes,
not merely an upper consistency bound or an empirical coincidence.
No number of Turing-machine states is inferred from this bit count.

\subsection{The equations needed by the reduction}

The unchanged signature is
\[
 d^0,N^2,V^2,P^2,A^2,Q^1,W^2,
 \mathsf{Target}^1,\mathsf{Next}^2,\mathsf{Mem}^2,
 p_E^2,p_\varepsilon^2.
\]
The first equation of $\mathcal E_{713}$ is the published Sheffer single
axiom
\[
 N(N(x,N(N(y,x),x)),N(y,N(z,x)))=y.
 \tag{Sh1}
\]
It makes the reduct a Boolean algebra with $N$ interpreted as NAND.
The primary proof source is McCune, Veroff, Fitelson, Harris, Feist and Wos,
\emph{Short Single Axioms for Boolean Algebra}, Journal of Automated
Reasoning 29 (2002), 1--16; the authors' supplementary site includes
the explicit Otter and Prover9 derivations:
\url{https://www.cs.unm.edu/~veroff/BA/}.

The equation to be removed is
\[
                    d=N(d,N(d,d)).
                    \tag{T}
\]
In a Boolean algebra its right side is $d\mathbin{\mathrm{NAND}}\neg d=1$,
so $(T)$ states exactly $d=1$. Throughout the proof, $1$ and $0$ denote
the \emph{derived Boolean} top and bottom. No prior identification of the
native symbol $d$ with either of them is assumed.

Two retained equations are
\begin{align*}
 V(P(x,s),d)&=x,\tag{Head}\\
 p_\varepsilon(s,P(a,b))
    &=N(p_\varepsilon(s,a),p_\varepsilon(s,b)).\tag{NAND class}
\end{align*}
For completeness, the other unchanged native and incidence equations are
\begin{align*}
 V(P(x,s),A(j,d))&=V(s,j),\tag{Native tail}\\
 p_\varepsilon(s,A(i,j))&=p_E(V(s,i),V(s,j)),\tag{Atom class}\\
 p_\varepsilon(s,Q(a))
   &=N\bigl(p_\varepsilon(P(W(s,a),s),a),
            p_\varepsilon(P(z,s),a)\bigr).\tag{Binder}
\end{align*}
Together with Sh1, Head, NAND class and the closure equation below,
these are all seven equations of $\mathcal E_{688}$. All variables are
universally quantified. The binder denotes existential negation, with
$W$ selecting a witness; its inherited calibration is explained in
\texttt{reference/kernel\_section.tex}.

The remaining equation needed for the derivation is the fully expanded
closure equation. To state its Boolean structure, use nine arguments
$d,r,e,f,a,b_0,c,u,v$ and define the following \emph{Boolean terms}:
\begin{align*}
 L&=N(d,N(c,N(a,N(d,b_0)))),\\
 B^-&=N(r,N(d,N(e,N(d,f)))),\\
 H^-&=N(N(L,N(d,u)),N(u,N(d,v))),\\
 F(d,r,e,f,a,b_0,c,u,v)&=N(B^-,H^-).
\end{align*}
The last equation in the literal $713$-bit packet is $d=F$, after
substitution of the following native terms. The descriptor variable is
$b$; the Boolean argument $b_0$ is a different notation.
\begin{align*}
 h&=\mathsf{Target}(b),&
 n_z&=\mathsf{Next}(b,z),&
 m_z&=\mathsf{Mem}(b,z),\\
 R_b(x,y)&=p_\varepsilon(P(y,P(x,b)),V(b,d)),&&&
\end{align*}
and
\begin{align*}
 r&=R_b(h,n_z),&e&=p_E(m_z,n_z),&f&=p_E(m_z,z),\\
 a&=p_E(W(x,d),x),&b_0&=p_E(x,h),&c&=R_b(x,y),\\
 u&=p_E(y,h),&v&=R_b(\mathsf{Next}(b,y),y).&&
\end{align*}
All the displayed abbreviations are expanded in the delivered AST and
equation listing. The new checker reconstructs this entire equation
independently and compares it literally with the inherited source.

\subsection{A Boolean cofactor forces truth}

\paragraph{Lemma.}
In every Boolean algebra, with no assumption on $d$,
\[
 \neg d\wedge F(d,r,e,f,a,b_0,c,u,v)=\neg d\wedge r.
 \tag{Cofactor}
\]
Consequently, $r=1$ and $d=F(d,r,e,f,a,b_0,c,u,v)$ imply $d=1$.

\paragraph{Proof.}
First set the formal Boolean argument $d$ to $0$. Every term $N(d,t)$
then equals $1$, giving
\[
 L=1,\qquad B^-=\neg r,\qquad
 H^-=N(N(1,1),N(u,1))=N(0,\neg u)=1.
\]
It follows that $F(0,r,e,f,a,b_0,c,u,v)=N(\neg r,1)=r$.
The Boolean cofactor identity now gives the displayed equality.

For an explicit all-algebra justification, let $q=\neg d$ and restrict
Boolean values to the principal interval $[0,q]$. The map
$t\mapsto t\wedge q$ preserves the Boolean operations when the interval
uses its relative complement and its own top $q$. It sends $d$ to $0$,
so the preceding calculation in the interval proves $(\mathrm{Cofactor})$.
This is a calculation on already evaluated predicate values. It does not
project the native domain or assert that $P,V,W$ commute with the map.
If $r=1$ and $d=F$, the identity yields
$0=\neg d\wedge d=\neg d$, hence $d=1$.
\hfill$\square$

\subsection{A constant true class supplies the required instance}

Define the concrete ground terms
\[
              k=P(d,P(d,d)),\qquad b_*=P(k,d).
\]
The NAND-class equation gives, for every native stream argument $s$,
\begin{align*}
 p_\varepsilon(s,k)
 &=N\bigl(p_\varepsilon(s,d),
          N(p_\varepsilon(s,d),p_\varepsilon(s,d))\bigr)\\
 &=1.
\end{align*}
This uses the Boolean tautology $t\mathbin{\mathrm{NAND}}\neg t=1$.
It makes no assumption about the class denoted by $d$ or the value
$p_\varepsilon(s,d)$. In particular it does not use the equation being
deleted.

The native head equation gives $V(b_*,d)=k$. Substitute $b=b_*$ and
$x=y=z=d$ in the retained closure equation. Its first matrix value is
therefore
\[
 r=p_\varepsilon\!\left(
  P(\mathsf{Next}(b_*,d),P(\mathsf{Target}(b_*),b_*)),k\right)=1.
\]
The cofactor lemma applied to this single ground instance forces $d=1$.
Thus the missing equation $(T)$ is a consequence of the retained
equations.

\paragraph{Theorem.}
Let $\mathcal E_{688}$ be the seven-equation theory obtained by deleting
$(T)$ from $\mathcal E_{713}$. Every $\mathcal E_{688}$ algebra satisfies
$(T)$. Hence $\mathcal E_{688}$ and $\mathcal E_{713}$ have exactly the
same models and the same equational consequences. In particular, the
ZFC-formalizable inherited calibration gives
\[
 \mathcal E_{688}\not\vdash N(d,d)=d
 \quad\longleftrightarrow\quad
 \mathcal E_{713}\not\vdash N(d,d)=d
 \quad\longleftrightarrow\quad
 \operatorname{Con}(\mathrm{ZFC}).
\]

\paragraph{Proof.}
One model inclusion follows by deleting an axiom. The opposite inclusion
is the preceding derivation. Equational soundness and completeness give
equality of consequences; alternatively one can insert the derived
equation whenever the old theory uses $(T)$. The final consistency
comparison is precisely the ordinary $\mathcal E_{713}$/raw-EPSC/ZFC
calibration in the incoming archive, with its stated dependencies.
The present reduction introduces no new set-theoretic assumption.
\hfill$\square$

The argument is ordinary mathematics formalizable in ZFC. A fully
expanded contextual rewrite trace deriving $(T)$ directly from Sh1 has
not been emitted. The supplied ground-instance recipe and Boolean
identity are an explicit proof recipe, not a claimed independent
proof-assistant endpoint or an internal PA proof-code transformer.

\subsection{Exact accounting and reproducibility}

The ordered signature, its arities, the variable capacity, the opcode
width, both endpoints, and every other equation are unchanged. The
original explicit truth equation has six tree nodes across both sides.
The codec uses one leading equation bit and four bits per node, so its
deletion saves exactly
\[
                        1+4\cdot6=25\text{ bits}.
\]
The new equation-node counts are
$14,6,10,12,12,17,82$, totaling $153$. The complete packet count is
\[
 \underbrace{8+3\cdot12+8}_{52\text{ header}}
 +\underbrace{7+1+4\cdot153}_{620\text{ equations}}
 +\underbrace{4\cdot4}_{16\text{ endpoints}}
 =688.
\]

\begin{center}
\begin{tabular}{lrrrr}
\hline
Packet & Equations & Equation nodes & Header / endpoints & Total bits\\
\hline
Inherited $\mathcal E_{713}$ & 8 & 159 & 52 / 16 & 713\\
Derived-truth $\mathcal E_{688}$ & 7 & 153 & 52 / 16 & 688\\
\hline
\end{tabular}
\end{center}

The exact $688$ bit characters, excluding the final newline, have SHA-256
\begin{center}\small\ttfamily
9518c474861557da085c76ef77dd4141eecfa0befaec66d90acc81ac39264d97.
\end{center}

Run \texttt{python3 -B build\_e688.py} in the supplied
\texttt{formalisms} directory. Its independent AST/codec implementation
first reconstructs the exact inherited $713$ bits and checks their pin.
It then deletes only the displayed truth axiom, validates the resulting
arities and ground endpoints, round-trips the complete new packet, and
rejects every one of its $688$ strict truncations and three malformed
controls. The Boolean cofactor is checked on all $512$ assignments to its
nine independent Boolean arguments. These finite checks substantiate
transcription and the finite algebraic identity; the written proof is
what establishes the arbitrary-algebra redundancy.

The deliverable includes \texttt{E688.bits}, \texttt{E688.json}, the fully
expanded \texttt{E688.equations}, a ground-substitution recipe, and a
verification receipt. A fair ground-rewrite search can use this packet
with exactly the same abstract checker as $\mathcal E_{713}$. Its binary
implementation, full writer and initialization cost still require a
separate compilation. This reduction supplies a smaller calibrated
equational input; it is not itself a sub-$237$-state Turing machine.
